% !TEX root = ../main.tex
%% Chapitre 12 — Combinaisons linéaires et indépendance
%% Sources : frontend/cours/resume/terminal/resume_combinaisons_lineaires_terminal.txt
%%           frontend/exercices/terminal_spé_Maths/exercice_combinaisons_lineaires.txt
%% Tous les calculs ont été vérifiés par programme (calcul exact).
%% Exercices 10 et 11 : exercices inédits rédigés pour le livre.

\begin{ouverture}[bases]{12}{Combinaisons linéaires et indépendance}{Partie 2 · Géométrie dans l'espace}
  {Combinaisons linéaires\enspace·\enspace Indépendance\enspace·\enspace Bases\enspace·\enspace Repères\enspace·\enspace
   Coordonnées}
  \objectif{Écrire un vecteur comme combinaison linéaire d'autres vecteurs en résolvant un système.}
  \objectif{Reconnaître une famille libre ou liée de deux ou trois vecteurs.}
  \objectif{Démontrer que trois vecteurs forment une base de l'espace et décomposer un vecteur dans cette base.}
  \objectif{Utiliser les coordonnées dans un repère quelconque de l'espace.}
  \prerequis{Vecteurs, colinéarité et coplanarité (chapitre 11).}
  \prerequis{Systèmes d'équations linéaires ; déterminant de deux vecteurs du plan (Seconde).}
\end{ouverture}
\miseEnPageSouple

%% =====================================================================
%%  L'ESSENTIEL
%% =====================================================================
\essentiel

\essSection{1}{Familles libres et familles liées}

\begin{definition}[Indépendance linéaire]
Des vecteurs $\vec{u_1}$, …, $\vec{u_n}$ sont \textbf{linéairement indépendants} (on dit que la famille est \textbf{libre}) si la
seule combinaison linéaire nulle est celle dont tous les coefficients sont nuls :
\[a_1\vec{u_1}+\dots+a_n\vec{u_n}=\vec0\ \Longrightarrow\ a_1=\dots=a_n=0.\]
Sinon, la famille est \textbf{liée} : l'un des vecteurs est une combinaison linéaire des autres.
\end{definition}

\noindent
\begin{minipage}[t]{0.5\linewidth}\vspace{0pt}
{\renewcommand{\arraystretch}{1.35}\arrayrulecolor{otRule}%
\begin{tabularx}{\linewidth}{@{}lY@{}}
\toprule
\otTete{Famille} & \otTete{Libre si et seulement si} \\
\midrule
\otNom{deux vecteurs} & ils ne sont pas colinéaires \\
\otNom{trois vecteurs} & ils ne sont pas coplanaires \\
\otNom{quatre vecteurs} & jamais : dans l'espace, elle est toujours liée \\
\bottomrule
\end{tabularx}}
\end{minipage}\hfill
\begin{minipage}[t]{0.46\linewidth}\vspace{0pt}
\begin{propriete}[Dans le plan]
\raggedright
$\vec u\,(x\pv y)$ et $\vec v\,(x'\pv y')$ sont colinéaires si et seulement si leur \textbf{déterminant} est nul :
\[\det(\vec u,\vec v)=xy'-x'y=0.\]
\end{propriete}
\end{minipage}

\essSection{2}{Bases et repères de l'espace}

\noindent
\begin{minipage}[c]{0.56\linewidth}
\begin{theoreme}[Décomposition dans une base]
\raggedright
Si $\vec i$, $\vec j$, $\vec k$ ne sont pas coplanaires, ils forment une \textbf{base} de l'espace : tout vecteur $\vec u$ s'écrit de
façon \textbf{unique}
\[\vec u=x\vec i+y\vec j+z\vec k.\]
Le triplet $(x\pv y\pv z)$ est formé des \textbf{coordonnées} de $\vec u$ dans la base $(\vec i,\vec j,\vec k)$.
\end{theoreme}
\end{minipage}\hfill
\begin{minipage}[c]{0.4\linewidth}
\centering
\begin{tikzpicture}[otEspace,scale=1.15]
  \coordinate (O) at (0,0,0); \coordinate (M) at (2,1.4,1.5);
  \draw[otCache] (2,0,0) -- (2,1.4,0) -- (0,1.4,0);
  \draw[otCache] (2,1.4,0) -- (M);
  \draw[otCache] (2,0,0) -- (2,0,1.5) -- (M) -- (0,1.4,1.5) -- (0,0,1.5) -- (2,0,1.5);
  \draw[otCache] (0,1.4,0) -- (0,1.4,1.5);
  \draw[otVecN] (O) -- (1,0,0) node[below,font=\small] {$\vec i$};
  \draw[otVecN] (O) -- (0,1,0) node[above left,font=\small,inner sep=1pt] {$\vec j$};
  \draw[otVecN] (O) -- (0,0,1) node[left,font=\small] {$\vec k$};
  \draw[otVec] (O) -- (M);
  \node[otPt] at (O) {}; \node[otSommet,anchor=north east] at (O) {$O$};
  \node[otPtV] at (M) {}; \node[otSommet,anchor=south west] at (M) {$M$};
\end{tikzpicture}\\[0.5mm]
\legende{$\vect{OM}=2\vec i+1{,}4\,\vec j+1{,}5\,\vec k$.}
\end{minipage}

\begin{propriete}[Repère et coordonnées]
Un \textbf{repère} $(O\pv\vec i,\vec j,\vec k)$ est formé d'un point $O$ et d'une base. Les coordonnées d'un point $M$ sont celles du
vecteur $\vect{OM}$. Dans tout repère : $\vect{AB}$ a pour coordonnées $(x_B-x_A\pv y_B-y_A\pv z_B-z_A)$, le milieu de $[AB]$ a pour
coordonnées $\left(\frac{x_A+x_B}{2}\pv\frac{y_A+y_B}{2}\pv\frac{z_A+z_B}{2}\right)$, et les coordonnées de $\vec u+\vec v$ et de
$k\vec u$ se calculent coordonnée par coordonnée.
\end{propriete}

\begin{savoirfaire}[Montrer qu'une famille est une base, puis décomposer]
\textbf{Exemple.} Soit $\vec u\,(1\pv2\pv0)$, $\vec v\,(0\pv1\pv1)$, $\vec w\,(1\pv0\pv1)$.
\begin{itemize}[leftmargin=1.2em,itemsep=0.4mm,topsep=0.5mm]
  \item \textbf{Base :} si $a\vec u+b\vec v+c\vec w=\vec0$, alors $a+c=0$, $2a+b=0$ et $b+c=0$, d'où $c=-a$, $b=-2a$, puis
    $-3a=0$ : $a=b=c=0$. La famille est libre : c'est une base.
  \item \textbf{Décomposition de} $\vec s\,(1\pv5\pv0)$ : on résout $a+c=1$, $2a+b=5$, $b+c=0$. On obtient $a=2$, $b=1$, $c=-1$, soit
    $\vec s=2\vec u+\vec v-\vec w$ : les coordonnées de $\vec s$ dans la base $(\vec u,\vec v,\vec w)$ sont $(2\pv1\pv-1)$.
\end{itemize}
\end{savoirfaire}

\begin{demonstration}[Unicité de la décomposition]
Supposons $\vec u=x\vec i+y\vec j+z\vec k=x'\vec i+y'\vec j+z'\vec k$. Par différence,
$(x-x')\vec i+(y-y')\vec j+(z-z')\vec k=\vec0$. La famille $(\vec i,\vec j,\vec k)$ étant libre, tous les coefficients sont nuls :
$x=x'$, $y=y'$ et $z=z'$.
\end{demonstration}

\begin{exemple}[Deux repères pour un même cube]
\begin{codecote}[0.62]
Dans le cube $ABCDEFGH$, les vecteurs $\vect{AB}$, $\vect{AD}$, $\vect{AE}$ ne sont pas coplanaires : $\big(A\pv\vect{AB},\vect{AD},\vect{AE}\big)$
est un repère, dans lequel $G$ a pour coordonnées $(1\pv1\pv1)$ et $H$ a pour coordonnées $(0\pv1\pv1)$.\\[1mm]
Dans le repère $\big(B\pv\vect{BA},\vect{BC},\vect{BF}\big)$, on a $\vect{BG}=0\,\vect{BA}+1\,\vect{BC}+1\,\vect{BF}$ :
$G$ a pour coordonnées $(0\pv1\pv1)$, et $H$ a pour coordonnées $(1\pv1\pv1)$. Les coordonnées d'un point dépendent du repère.
\cote
\centering
\begin{tikzpicture}[otEspace,scale=2.2]
  \otCube
  \draw[otVecN] (A) -- (B); \draw[otVecN] (A) -- (D); \draw[otVecN] (A) -- (E);
  \draw[otVec] (B) -- (A); \draw[otVec] (B) -- (C); \draw[otVec] (B) -- (F);
\end{tikzpicture}
\end{codecote}
\end{exemple}

\begin{attention}[Pièges fréquents]
\begin{itemize}[label={\color{otOrange}\textbullet},leftmargin=1.4em,itemsep=0.6mm,topsep=0.5mm]
  \item Trois vecteurs deux à deux non colinéaires peuvent être coplanaires : il faut étudier les trois ensemble.
  \item Les coordonnées d'un vecteur dépendent de la base choisie, et de l'ordre des vecteurs de la base.
  \item Dans un repère quelconque, la longueur d'un vecteur ne se calcule pas par $\sqrt{x^2+y^2+z^2}$ : il faut un repère
    orthonormé (chapitre 13).
  \item Une famille contenant le vecteur nul est toujours liée.
\end{itemize}
\end{attention}

%% =====================================================================
%%  EXERCICES
%% =====================================================================
\exercices
\rubrique{Découverte}

\begin{exercice}[1]{Rappels : calculs dans le plan}
Dans un repère du plan, on considère $\vec u\,(3\pv2)$, $\vec v\,(-1\pv4)$ et $\vec w\,(5\pv-2)$.
\begin{questions}
  \item Calculer les coordonnées de $\vec u+\vec v$, $2\vec u-\vec w$, $\vec u+\vec v+\vec w$ et $3\vec u-2\vec v+\vec w$.
  \item Déterminer les coordonnées du vecteur $\vect{AB}$, avec $A(2\pv5)$ et $B(7\pv1)$.
  \item Calculer $\det\big((2\pv5),(3\pv7)\big)$ et $\det\big((6\pv9),(4\pv6)\big)$. Que peut-on en conclure ?
  \item Pour quelle valeur de $m$ les vecteurs $\vec p\,(m\pv2)$ et $\vec q\,(3\pv6)$ sont-ils colinéaires ?
\end{questions}
\end{exercice}

\begin{exercice}[1]{Colinéarité et alignement}
\begin{questions}
  \item Les vecteurs $\vec u\,(4\pv6)$ et $\vec v\,(2\pv3)$ sont-ils colinéaires ? Et $\vec a\,(3\pv-2)$ et $\vec b\,(6\pv-5)$ ?
  \item Déterminer $k$ pour que $\vec p\,(k\pv8)$ et $\vec q\,(3\pv4)$ soient colinéaires.
  \item Les points $A(1\pv2)$, $B(4\pv5)$ et $C(7\pv8)$ sont-ils alignés ?
  \item Les vecteurs $\vec u\,(-2\pv3)$ et $\vec v\,(4\pv-6)$ forment-ils une base du plan ? Et $\vec m\,(5\pv2)$ et $\vec n\,(3\pv7)$ ?
\end{questions}
\end{exercice}

\begin{exercice}[1]{Décomposer un vecteur}
Dans un repère du plan, on considère $\vec u\,(1\pv2)$ et $\vec v\,(3\pv-1)$.
\begin{questions}
  \item Vérifier que $\vec u$ et $\vec v$ forment une base du plan.
  \item Exprimer $\vec w\,(7\pv1)$, puis $\vec a\,(5\pv8)$, comme combinaisons linéaires de $\vec u$ et $\vec v$.
  \item Le vecteur $\vec b\,(2\pv4)$ peut-il s'exprimer à l'aide de $\vec u$ seulement ?
  \item Déterminer les coordonnées des vecteurs $a\vec u+b\vec v$ tels que $a+b=5$, en fonction de $a$.
\end{questions}
\end{exercice}

\rubrique{Entraînement}

\begin{exercice}[2]{Changer de base dans le plan}
On considère $\vec{e_1}\,(1\pv0)$, $\vec{e_2}\,(0\pv1)$, $\vec u\,(2\pv3)$ et $\vec v\,(1\pv-1)$.
\begin{questions}
  \item Exprimer $\vec u$ dans la base $(\vec{e_1},\vec{e_2})$.
  \item Montrer que $(\vec u,\vec v)$ est une base du plan.
  \item Exprimer $\vec{e_1}$ puis $\vec{e_2}$ dans la base $(\vec u,\vec v)$.
  \item Déterminer les réels $a$ et $b$ tels que $a\vec u+b\vec v=(5\pv4)$. Retrouver ce résultat à l'aide de la question 3.
\end{questions}
\end{exercice}

\begin{exercice}[2]{Une famille liée}
On considère les vecteurs du plan $\vec{u_1}\,(1\pv1)$, $\vec{u_2}\,(2\pv-1)$ et $\vec{u_3}\,(5\pv1)$.
\begin{questions}
  \item Montrer que $\vec{u_1}$ et $\vec{u_2}$ sont linéairement indépendants.
  \item Exprimer $\vec{u_3}$ comme combinaison linéaire de $\vec{u_1}$ et $\vec{u_2}$, puis écrire une relation
    $a\vec{u_1}+b\vec{u_2}+c\vec{u_3}=\vec0$ à coefficients entiers non tous nuls.
  \item Expliquer pourquoi trois vecteurs du plan forment toujours une famille liée.
  \item Le vecteur $\vec w\,(0\pv5)$ peut-il s'exprimer comme combinaison linéaire de $\vec{u_1}$ et $\vec{u_2}$ ?
\end{questions}
\end{exercice}

\begin{exercice}[2]{Applications géométriques}
\begin{codecote}[0.56]
Dans un repère du plan, on considère les points $A(2\pv1)$, $B(5\pv3)$, $C(1\pv4)$ et $D(4\pv6)$.
\begin{questions}
  \item Calculer les coordonnées du milieu $M$ de $[AB]$.
  \item Les droites $(AB)$ et $(CD)$ sont-elles parallèles ?
  \item Déterminer les coordonnées du point $E$ tel que $ABCE$ soit un parallélogramme.
  \item Exprimer $\vect{AD}$ comme combinaison linéaire de $\vect{AB}$ et $\vect{AC}$. Interpréter.
  \item Les points $A$, $B$ et $D$ sont-ils alignés ?
\end{questions}
\cote
\centering
\begin{tikzpicture}
\begin{axis}[otGraphe, width=62mm, xmin=-3, xmax=6.5, ymin=-0.5, ymax=7, xtick={-2,...,6}, ytick={1,...,6},
  grid=major, grid style={otRule!60,line width=0.3pt}]
  \addplot[only marks,mark=*,mark size=1.6pt,otNavy] coordinates {(2,1) (5,3) (1,4) (4,6)};
  \node[font=\small,anchor=north] at (axis cs:2,1) {$A$};
  \node[font=\small,anchor=west] at (axis cs:5,3) {$B$};
  \node[font=\small,anchor=east] at (axis cs:1,4) {$C$};
  \node[font=\small,anchor=south] at (axis cs:4,6) {$D$};
\end{axis}
\end{tikzpicture}
\end{codecote}
\end{exercice}

\begin{exercice}[2]{Équations vectorielles}
On considère les vecteurs du plan $\vec u\,(2\pv1)$, $\vec v\,(-1\pv3)$ et $\vec w\,(3\pv7)$.
\begin{questions}
  \item Résoudre l'équation vectorielle $x\vec u+y\vec v=\vec w$.
  \item Le vecteur $\vec a\,(1\pv5)$ est-il une combinaison linéaire de $\vec u$ et $\vec v$ ?
  \item Trouver les réels $a$ et $b$ tels que $a\vec u+b\vec v=(0\pv10)$.
  \item Exprimer $\vec v$ en fonction de $\vec u$ et de $\vec w$.
  \item Les points $M(2\pv1)$, $N(-1\pv3)$ et $P(8\pv-3)$ sont-ils alignés ?
\end{questions}
\end{exercice}

\rubrique{Problèmes}

\begin{exercice}[2]{Vecteurs de l'espace}
Dans un repère de l'espace, on considère $\vec u\,(1\pv0\pv2)$, $\vec v\,(0\pv1\pv1)$ et $\vec w\,(2\pv1\pv5)$.
\begin{questions}
  \item Exprimer $\vec w$ comme combinaison linéaire de $\vec u$ et $\vec v$. La famille $(\vec u,\vec v,\vec w)$ est-elle libre ?
  \item On pose $\vec t\,(0\pv0\pv1)$. Montrer que $(\vec u,\vec v,\vec t)$ est une base de l'espace.
  \item Les points $O$, $A(1\pv0\pv2)$, $B(1\pv1\pv3)$ et $C(3\pv2\pv7)$ sont-ils coplanaires ?
\end{questions}
\end{exercice}

\begin{exercice}[3]{Coplanarité et base dans l'espace}
On considère $\vec u\,(1\pv2\pv1)$, $\vec v\,(2\pv1\pv3)$, $\vec w\,(3\pv3\pv4)$ et $\vec t\,(1\pv0\pv0)$.
\begin{questions}
  \item Montrer que $\vec u$ et $\vec v$ ne sont pas colinéaires, puis exprimer $\vec w$ en fonction de $\vec u$ et $\vec v$.
  \item Les vecteurs $\vec u$, $\vec v$ et $\vec w$ sont-ils coplanaires ? La famille $(\vec u,\vec v,\vec w)$ est-elle une base ?
  \item Montrer que $(\vec u,\vec v,\vec t)$ est une base de l'espace.
  \item Les points $O$, $A(1\pv2\pv1)$, $B(3\pv3\pv4)$ et $C(5\pv4\pv7)$ sont-ils coplanaires ?
\end{questions}
\end{exercice}

\begin{exercice}[3]{Coordonnées dans une base non orthogonale}
Dans un repère $(O\pv\vec i,\vec j,\vec k)$, on considère $\vec u\,(1\pv1\pv0)$, $\vec v\,(0\pv1\pv1)$ et $\vec w\,(1\pv0\pv1)$.
\begin{questions}
  \item Démontrer que $(\vec u,\vec v,\vec w)$ est une base de l'espace.
  \item Déterminer les coordonnées des vecteurs $\vec s\,(4\pv3\pv1)$ et $\vec r\,(1\pv1\pv1)$ dans la base $(\vec u,\vec v,\vec w)$.
  \item Le point $M$ a pour coordonnées $(2\pv-1\pv3)$ dans le repère $(O\pv\vec u,\vec v,\vec w)$. Quelles sont ses coordonnées dans
    le repère $(O\pv\vec i,\vec j,\vec k)$ ?
  \item Plus généralement, exprimer les coordonnées $(x\pv y\pv z)$ d'un point dans $(O\pv\vec i,\vec j,\vec k)$ en fonction de ses
    coordonnées $(X\pv Y\pv Z)$ dans $(O\pv\vec u,\vec v,\vec w)$.
\end{questions}
\end{exercice}

\begin{exercice}[3]{Une famille dépendant d'un paramètre}
Pour tout réel $m$, on considère les vecteurs $\vec u\,(1\pv m\pv1)$, $\vec v\,(m\pv1\pv1)$ et $\vec w\,(1\pv1\pv m)$.
\begin{questions}
  \item Que dire de la famille $(\vec u,\vec v,\vec w)$ lorsque $m=1$ ?
  \item Calculer $\vec u+\vec v+\vec w$ lorsque $m=-2$. Conclure.
  \item On suppose $m\neq1$. Soit $a$, $b$, $c$ des réels tels que $a\vec u+b\vec v+c\vec w=\vec0$. En soustrayant deux à deux les
    équations obtenues, montrer que $a=b=c$, puis que $(m+2)a=0$.
  \item En déduire les valeurs de $m$ pour lesquelles $(\vec u,\vec v,\vec w)$ est une base de l'espace.
  \item Pour $m=0$, exprimer le vecteur $\vec s\,(2\pv2\pv2)$ dans la base $(\vec u,\vec v,\vec w)$.
\end{questions}
\end{exercice}

\begin{exercice}[3]{Centre de gravité d'un tétraèdre}
\begin{codecote}[0.58]
$ABCD$ est un tétraèdre. On travaille dans le repère $\big(A\pv\vect{AB},\vect{AC},\vect{AD}\big)$. On note $G$ le centre de gravité
du triangle $BCD$, défini par $\vect{GB}+\vect{GC}+\vect{GD}=\vec0$, et $O$ le point tel que
$\vect{OA}+\vect{OB}+\vect{OC}+\vect{OD}=\vec0$.
\begin{questions}
  \item Donner les coordonnées de $A$, $B$, $C$, $D$, puis montrer que $G\left(\frac13\pv\frac13\pv\frac13\right)$.
  \item Montrer que $O\left(\frac14\pv\frac14\pv\frac14\right)$, puis que $A$, $O$ et $G$ sont alignés. Préciser la position de $O$ sur
    $[AG]$.
  \item $I$ et $J$ sont les milieux de $[AB]$ et $[CD]$. Montrer que $O$ est le milieu de $[IJ]$.
  \item En déduire que les trois segments joignant les milieux de deux arêtes opposées sont concourants.
\end{questions}
\cote
\centering
\begin{tikzpicture}[scale=1.3]
  \coordinate (A) at (0,0); \coordinate (B) at (1.6,-0.9); \coordinate (C) at (3,0.15); \coordinate (D) at (1.35,2.35);
  \coordinate (G) at ($1/3*(B)+1/3*(C)+1/3*(D)$); \coordinate (O) at ($1/4*(A)+1/4*(B)+1/4*(C)+1/4*(D)$);
  \coordinate (I) at ($(A)!0.5!(B)$); \coordinate (J) at ($(C)!0.5!(D)$);
  \fill[otGreen,opacity=0.12] (B) -- (C) -- (D) -- cycle;
  \draw[otVu] (A) -- (B) -- (C) -- (D) -- cycle; \draw[otVu] (B) -- (D);
  \draw[otCache] (A) -- (C);
  \draw[otConstr] (A) -- (G);
  \draw[otGreenDark,line width=0.8pt,dash pattern=on 2.2pt off 1.6pt] (I) -- (J);
  \node[otPtV] at (G) {}; \node[otSommet,anchor=west] at (G) {$G$};
  \node[otPt] at (O) {}; \node[otSommet,anchor=north] at (O) {$O$};
  \foreach \p/\a in {I/north east,J/west}{\node[otPtV] at (\p) {}; \node[otSommet,anchor=\a] at (\p) {$\p$};}
  \node[otSommet,anchor=east] at (A) {$A$}; \node[otSommet,anchor=north] at (B) {$B$};
  \node[otSommet,anchor=west] at (C) {$C$}; \node[otSommet,anchor=south] at (D) {$D$};
\end{tikzpicture}
\end{codecote}
\end{exercice}
